\documentclass[11pt]{article}
\usepackage[a4paper,margin=28mm]{geometry}
\usepackage{amsmath,amssymb,amsthm,mathtools,hyperref}
\newtheorem{theorem}{Theorem}
\newtheorem{lemma}[theorem]{Lemma}
\newtheorem{corollary}[theorem]{Corollary}
\newtheorem{proposition}[theorem]{Proposition}
\theoremstyle{remark}\newtheorem{remark}[theorem]{Remark}
\newcommand{\M}{\mathcal M^1}
\newcommand{\Pp}{\Phi_{S,r}}
\newcommand{\D}{\mathrm D}
\newcommand{\supp}{\operatorname{supp}}
\title{Order-controlled intrinsic extensions and sufficient certificates\ for non-Archimedean Bedford--Taylor actions}
\author{Internal continuation after the R24 audit}
\date{14 November 2026}
\begin{document}
\maketitle
\begin{abstract}
We establish quantitative order continuity of the already constructed
continuous-coefficient curvature products, identify their order completions
with actual currents, and extend the products uniquely to these completions.
We also prove that controlled approximation by already annihilated
degree-one forms suffices for the quadratic action obstruction to vanish.
These are consequences of accepted internal inputs, not a solution of the
general analytic relation problem. In particular, no density assertion,
new principal descent theorem, or claim of literature priority is made.
\end{abstract}
\tableofcontents

\section{Scope and precise inputs}

Work over the original finite extension $K/\mathbf Q_p$, with its unchanged
absolute value, on an ordinary open $U$ of a good, Hausdorff, boundaryless,
quasi-smooth analytic space, locally of pure dimension $n$. All equalities
of currents are on the full original compactly supported smooth CLD test
sheaf. All actual moment refinements are retained. Write $D=d'd''$ with
$d'd''=-d''d'$. An allowed background $S$ has bidegree $(p,p)$ and is
symmetric, middle-positive and separately closed. No stronger positive cone
or geometric presentation of $S$ is imposed.

The input sheaf $\M$ consists of symmetric coefficient-Radon $(1,1)$
currents locally admitting finite presentations
\[
 A=\sum_i f_iD[g_i],\qquad f_i\in C(U,\mathbf R),
\]
where $g_i$ are locally bounded original classical psh functions; signed
differences are handled by finite expansion. Equality is actual current
equality, not equality of formal expressions. Its positive cone is the
original middle cone. Frozen R6 supplies symmetric, positive,
$C(U,\mathbf R)$-multilinear products
\[
 \Phi_{S,r}:(\M)^r\longrightarrow\mathcal D'^{p+r,p+r},\qquad
 \Phi_{S,r}(A_1,\ldots,A_r)
 =\sum_{i_1,\ldots,i_r}\prod_j f_{j,i_j}
       T_{S,r}(g_{1,i_1},\ldots,g_{r,i_r}),
 \tag{I1}
\]
for $r\le n-p$. Values are coefficient-Radon currents. This operation is
ordinary-local and respects all representations and actual refinements.
For $r=1$ put $M_S=\Phi_{S,1}$. The same notation agrees with the frozen
degree-one action on genuine symmetric analytic delta-forms. On
nonsymmetric genuine $B^{1,1}$ arguments in (I3), $M_S$ denotes the
broader frozen QA degree-one action. All order estimates and completions
in this note are restricted to symmetric arguments.

We use the frozen finite-frame coefficient domination: if $R$ is positive
and $T$ symmetric with $-cR\le T\le cR$, then, for each fixed original
compact test $\eta$, there is a finite positive control
\[
 |\langle T,\eta\rangle|\le c\,N_{R,\eta}<\infty.       \tag{I2}
\]
On each finite moment presentation and compact stage this also controls
the coefficient measures. It follows by applying positivity to $cR\pm T$
and the fixed finite positive controls for $\eta$. These are the original
middle-cone controls, not an assertion about a stronger cone. In particular,
an order-bounded limit is coefficient-Radon, the cone is weakly closed,
and differential operators are continuous for distributional convergence.

Frozen QA supplies, for actual analytic $\beta\in B^{0,1}$ and
$\zeta\in B^{1,0}$,
\[
 M_S(d'\beta)=d'(\beta\wedge S),\qquad
 M_S(d''\zeta)=d''(\zeta\wedge S).                       \tag{I3}
\]
These are full derivatives with every polyhedral boundary term present.
Consequently $DM_S(d'\beta+d''\zeta)=0$.

The exact source files and hashes are in the accompanying dependency ledger.
These inputs are adopted in their frozen scope; the present note does not
independently reprove the entire R1--R24 chain.

\section{A sharp order estimate}

\begin{lemma}[Signed products controlled by positive products]\label{lem:order}
Suppose $Q_i\in\M$ are positive and
$-M_iQ_i\le A_i\le M_iQ_i$ for $M_i\ge0$. Then
\[
 -\Bigl(\prod_iM_i\Bigr)R_S\le
 \Phi_{S,r}(A_1,\ldots,A_r)\le
 \Bigl(\prod_iM_i\Bigr)R_S,
 \qquad R_S=\Phi_{S,r}(Q_1,\ldots,Q_r)\ge0.             \tag{OC1}
\]
If $-M_iQ_i\le A_i,B_i\le M_iQ_i$ and
$-\delta_iQ_i\le A_i-B_i\le\delta_iQ_i$, then
\[
 -E R_S\le \Phi_{S,r}(A_1,\ldots,A_r)
               -\Phi_{S,r}(B_1,\ldots,B_r)\le E R_S,
 \qquad E=\sum_i\delta_i\prod_{j\ne i}M_j.             \tag{OC2}
\]
\end{lemma}
\begin{proof}
Put $A_i^\pm=(M_iQ_i\pm A_i)/2$. Each is a positive member of $\M$,
$A_i=A_i^+-A_i^-$ and $A_i^++A_i^-=M_iQ_i$. Expanding the product by
multilinearity gives a signed sum of positive currents. Its positive and
negative subsums are both bounded by the sum of all terms, which is
$\Phi_{S,r}(M_1Q_1,\ldots,M_rQ_r)$. More exactly, the unsigned sum
plus or minus the signed sum is twice one of those positive subsums.
This proves (OC1), including zero values of $M_i$. Telescope the difference
of the two products one slot at a time and apply (OC1) to each summand.
Adding the resulting order inequalities proves (OC2). No positivity of
a signed tail is assumed.
\end{proof}

The estimate holds for every allowed $S$ with the same numerical $E$.
The controlling mass $N_{R_S,\eta}$ depends on $S$ and the test; there is
no uniform bound on the masses of all backgrounds.

\section{Completion as actual currents}

For one positive $Q\in\M(U)$ set
\[
 \begin{aligned}
 E_Q(U)&=\{A\in\M(U): -cQ\le A\le cQ\text{ for some }c<\infty\},\\
 \|A\|_Q&=\inf\{c:-cQ\le A\le cQ\}.
 \end{aligned}                                         \tag{OC3}
\]
The case $Q=0$ gives the zero space. Every finite presentation belongs
locally to such an ideal: after bounding its coefficients, a sum of the
positive curvatures in that presentation is a control. No one global
control on an arbitrary open is asserted.

\begin{theorem}[Faithful realization of each order completion]\label{thm:completion}
The norm completion $\widehat E_Q(U)$ embeds isometrically into actual
symmetric coefficient-Radon $(1,1)$ currents on $U$. Its image is exactly
the closure of $E_Q(U)$ in the order norm determined by $Q$ in the ambient
space of such currents. If the same actual current is obtained with
controls $Q$ and $R$, the two realizations agree in the completion for
$Q+R$.
\end{theorem}
\begin{proof}
The gauge is a norm: if its value is zero, then $\pm A\le\epsilon Q$
for every $\epsilon>0$; (I2) on every test implies $A=0$.
The triangle inequality follows by adding order bounds. Weak closedness
of the cone also shows that the infimum defining a finite norm is attained.

If $(A_j)$ is order-Cauchy, (I2) makes its values Cauchy on every test.
The local finite-frame bound makes the limit a distribution and, on each
compact presentation, coefficient-Radon. Its symmetry and compatibility
between presentations pass to the limit. Write the resulting actual
current as $A$. If $\|A_j-A_k\|_Q\le\epsilon$ for $j,k\ge N$, let
$k\to\infty$ in the order inequalities to obtain
\[
                -\epsilon Q\le A_j-A\le\epsilon Q
                       \quad(j\ge N).                   \tag{OC4}
\]
Thus convergence is in the ambient order norm. If $A=0$, (OC4) says
the original sequence represents zero in the completion, proving
injectivity; (OC4) also proves the norm assertion and the characterization
of the image. Since $Q\le Q+R$ and $R\le Q+R$, both completion sequences
are Cauchy for the larger control. Their actual limits agree, so
injectivity for $Q+R$ identifies them.
\end{proof}

Define $\M_{\rm ord}$ locally: an actual symmetric coefficient-Radon
current belongs if every point has an ordinary neighborhood on which it
lies in one of these embedded completions. This is a sheaf by its local
definition and the sheaf property of currents. It is not a claim that a
single global order bound survives arbitrary gluing. Local addition uses
$Q+R$, and multiplication by continuous functions uses their local boundedness.

\begin{theorem}[Intrinsic extension of the mixed operations]\label{thm:extension}
The operations (I1) extend uniquely to symmetric locally
$C(U,\mathbf R)$-multilinear maps on $(\M_{\rm ord})^r$, continuous in
each common finite collection of local order controls. They are independent
of approximation sequences, controls and moment presentations, preserve
positivity, commute with ordinary restriction and gluing, and satisfy
(OC1)--(OC2). Their supports satisfy
\[
 \supp\Phi_{S,r}(A_1,\ldots,A_r)
       \subseteq\supp S\cap\bigcap_i\supp A_i.          \tag{OC5}
\]
If mixed products along one such approximation are separately closed,
their limit is separately closed. Separate closedness of the limit factors
alone is not claimed to imply this hypothesis.
\end{theorem}
\begin{proof}
Choose local controls $Q_i$ and Cauchy approximations $A_{i,j}$.
Their norms are bounded, so (OC2) and (I2) make the mixed products Cauchy
in coefficient measure on every fixed compact presentation. They have
a coefficient-Radon limit. Applying (OC2) to different approximations
proves uniqueness for these controls. For other controls $R_i$, replace
each by $Q_i+R_i$ and apply Theorem~\ref{thm:completion}; the same estimate
proves equality of the resulting limits. All algebraic properties follow
by taking limits of the frozen identities.

If $A_i\ge0$, (OC4) yields positive approximants
$A_{i,j}+\epsilon_{i,j}Q_i\ge0$ with $\epsilon_{i,j}\to0$.
They have the same limits. Positivity follows from (I1) and weak closedness
of the cone. Passing the order inequalities to the limit gives (OC1)--(OC2).
Restriction maps are contractive for the restricted controls. The local
constructions therefore agree on intersections and glue as ordinary
currents. On an open where one $A_i$ is zero, its completion class is zero
by Theorem~\ref{thm:completion}; multilinearity makes the product zero.
The same is true where $S=0$, proving (OC5). Distributional continuity
of $d'$ and $d''$ proves the last assertion.
\end{proof}

\begin{remark}[What the completion does not assert]
This is an intrinsic extension by a specified convergence mechanism.
We have not proved that it strictly enlarges $\M$, that all positive closed
currents belong to it, that products of arbitrary currents exist, or that
all actual higher-degree analytic relations are killed. The theorem does
not alter the carrier-CLN seminorm or its relation ideal. Closedness of
actions of arbitrary closed degree-one delta-forms remains a separate issue.
\end{remark}

\section{Controlled exact approximation and the quadratic obstruction}

For a finite actual PL relation
\[
 P=\sum_\nu c_\nu[h_\nu][g_\nu],\qquad
 \sum_\nu c_\nu D_\delta h_\nu\wedge D_\delta g_\nu=0,
\]
QA defines the symmetric actual degree-one primitive
\[
 \alpha_P=\frac12\sum_\nu c_\nu
       (h_\nu D_\delta g_\nu+g_\nu D_\delta h_\nu),
 \quad D_\delta\alpha_P=0,
 \quad\mathcal A_S(P)=DM_S(\alpha_P).                    \tag{OC6}
\]
The operation on the right is already defined, without assuming a
general degree-two action on delta-forms.

\begin{theorem}[Controlled approximation suffices]\label{thm:density}
Let $\alpha\in B^{1,1}(U)$ be symmetric. Suppose actual symmetric
degree-one forms $a_j$, a fixed positive $Q\in\M(U)$ and
$\epsilon_j\downarrow0$ satisfy
\[
 -\epsilon_jQ\le[\alpha]-[a_j]\le\epsilon_jQ,
 \qquad DM_S(a_j)=0\quad\text{for every allowed }S.      \tag{OC7}
\]
The data are chosen on one representative before $S$ and the tests.
Then $DM_S(\alpha)=0$ for every allowed $S$, with
\[
 |\langle DM_S(\alpha),\eta\rangle|
             \le \epsilon_j N_{M_S(Q),D\eta}.           \tag{OC8}
\]
One sufficient, non-circular choice is
$a_j=d'\beta_j+d''\zeta_j$ in the full actual delta sheaf.
In particular, (OC7) for $\alpha=\alpha_P$ proves
$\mathcal A_S(P)=0$ for all backgrounds in legal degree.
\end{theorem}
\begin{proof}
R6 positive scalar-module transfer gives
$-\epsilon_jM_S(Q)\le M_S(\alpha)-M_S(a_j)\le
\epsilon_jM_S(Q)$. Apply (I2) to the fixed test $D\eta$.
Distributional duality, whose fixed convention sign disappears inside
absolute values, gives (OC8). The right side tends to zero for each $S$
and each test. No derivative convergence of coefficients is needed.
Formula (I3) proves the asserted sufficient exactness condition. Finally
use (OC6). The argument can also be applied on neighborhoods of each
point and then glued; it does not restrict a global full polar.
\end{proof}

\begin{corollary}[Fixed generators and no gauge condition on potentials]
Suppose on one representative
\[
 [\gamma]=\sum_{i=1}^N f_iD[g_i],\qquad
 A_j=\sum_{i=1}^N f_{i,j}D[g_i]=D[v_j],                 \tag{OC9}
\]
where the $g_i$ are finite continuous classical psh differences, all
coefficients are continuous, $f_{i,j}\to f_i$ locally uniformly, and
$v_j$ are finite continuous local classical psh differences. Then
$M_S(\gamma)$ is separately closed for every allowed $S$.
No uniform bound, normalization, or convergence of $v_j$ is required.
For genuine degree-one delta curvatures the local DC requirement for a
given continuous potential is also supplied by frozen AP.
\end{corollary}
\begin{proof}
Write $g_i=g_i^+-g_i^-$ and take the fixed positive represented control
$Q=\sum_i(D[g_i^+]+D[g_i^-])$. On a compactly contained ordinary
neighborhood, with $\epsilon_j=\max_i\|f_i-f_{i,j}\|_\infty$,
the error in (OC9) lies between $\pm\epsilon_jQ$.
Equivalently its action is a finite sum of uniformly small coefficients
times fixed coefficient-Radon currents $D(g_iS)$. Thus
$M_S(A_j)\to M_S(\gamma)$ on all tests. Scalar relation transfer gives
$M_S(A_j)=D(v_jS)$, which is separately closed. Distributional continuity
of both first derivatives proves the assertion. Locality handles the
whole representative.
\end{proof}

This corollary replaces the demand for one exact continuous potential
by a particular controlled approximation target. No such approximation
has been constructed for all RV/QT boundaries, and strict logical
weakening within that analytic class has not been established. RV4 and
the general AP boundary problem therefore remain open. A weak limit of
ambient currents alone does not satisfy the displayed order hypothesis.

\section{A directly verifiable first-order relation certificate}

\begin{proposition}[Scalar rows and actual support conormals]\label{prop:conormal}
Put $A_i=D[g_i]$ and $T_i=D(g_iS)$ for fixed finite continuous psh
differences $g_i$. Let $F_i=\supp A_i$ be the actual analytic support.
Then $\supp T_i\subset F_i$. If $h$ is original smooth and vanishes
on $F_i$, then $d'h\wedge T_i=d''h\wedge T_i=0$.

Suppose smooth covectors $\theta_i$ admit a finite identity on $U$
\[
 \theta_i=\sum_a f_{ai}\eta_a+\sum_j b_{ij}d'h_{ij},
 \qquad \sum_i f_{ai}A_i=0,
 \qquad h_{ij}|_{F_i}=0,                               \tag{SC1}
\]
with continuous scalar coefficients and original smooth $\eta_a,h_{ij}$.
Then $\sum_i\theta_i\wedge T_i=0$. The same assertion holds with
the columns exchanged. There is also the version in which $\theta_i$
and $\eta_a$ are actual PS covectors, $f_{ai},b_{ij}$ are continuous PS
scalars, and $h_{ij}$ remain original smooth, provided the displayed
identities are identities in the already established PS coefficient
module. No new product for arbitrary continuous covector coefficients
is inferred from that version.
\end{proposition}
\begin{proof}
On $U\setminus F_i$ the scalar curvature relation $A_i=0$ transfers
by R6 to $T_i=0$, proving the support assertion. Multiplication by $h$
gives $hT_i=0$. Since $T_i$ is separately closed, the smooth distributional
product rule gives the conormal assertion. Each scalar relation in (SC1)
transfers to $\sum_i f_{ai}T_i=0$. Multiply it by $\eta_a$, use the
conormal assertion for the remaining terms, and sum. All currents and
supports are the actual original ones. No finite skeleton replaces $F_i$.
\end{proof}

The proposition is a sufficient certificate, not a generation theorem.
One would still have to show that the particular covector relation in
AP24.9 is generated by certified scalar rows, support conormals and the
already settled ambient relations. Such a theorem has not been proved.
The simpler implication obtained by differentiating a scalar zero
$\sum_i f_iA_i=0$ does not apply merely from the closedness of
$\sum_i f_iA_i$.

\section{A necessary warning about approximation}

An order estimate is stronger than weak convergence. Precisely, let $C$
and $Q$ be positive coefficient-Radon currents and let a fixed positive
smooth complementary test form $\theta$ define local scalar measures
$\mu_T(\varphi)=\langle T,\varphi\theta\rangle$. Suppose on a compact
measurable set $E$ that $\mu_C(E)=c>0$ whereas $\mu_{A_j}(E)=0$ for
every proposed approximant. If
$-\epsilon_jQ\le C-A_j\le\epsilon_jQ$, then
\[
                       c\le\epsilon_j\mu_Q(E).          \tag{OC10}
\]
Local finiteness of $Q$ prevents $\epsilon_j\to0$. This is simply the
scalar-measure consequence of the order inequality. It applies, for
example, to real hyperplane delta mass and its diffuse smooth
regularizations. The real example is a convergence-topology diagnostic,
not a counterexample to KC12 or to analytic density in a correctly
chosen class. Any use for an original cap must separately verify the
displayed measure assumptions. Ordinary smoothing therefore does not
by itself supply the required controlled approximation.

\section{Source comparison and remaining task}

All proved assertions above are recorded as direct consequences or
functional-analytic extensions of frozen inputs. We found no matching
order-completion statement in the source modules inspected, but do not
claim an exhaustive priority audit. These estimates do not settle the
existence of an approximating family in (OC7) or (OC9).

For comparison, Mihatsch, \emph{On Tropical Intersection Theory},
arXiv:2107.12067v3, Lemma 4.5, gives a PL potential for a divisor on
$\mathbf R^n$. This is a statement about the ambient real space. We do
not replace a scalar cycle closed only on the full analytic quotient
by an ambient closed divisor. The different numbering in v2 is not used.

The next geometric objective is either to construct an intrinsic
controlled exact approximation for the actual quadratic primitive, or
to construct (OC9) for the actual RV scalar remainder, or to prove the
specific covector relation has a certificate of type (SC1) with any
additional already justified ambient generators. These are alternative
routes. A successful scalar-remainder construction still leaves the
balanced lift RV4 unless it is bypassed by a direct argument for
$\alpha_P$. KC12, OPEN-GQ, NC12, variable-PS and CTO retain their
frozen open status.
\end{document}
